% !TeX root = ../main.tex

% ==================================================
% D.4. 레이블과 위치 기반 선택
% ==================================================

\subsection{레이블과 위치 기반 선택}

\subsubsection{이전 학습과의 연결}

\zsh{D.2}에서는 Series, DataFrame, Index가 각각 독립적인 Python 객체이며, Index가 데이터의 레이블을 관리한다는 사실을 배웠다.

\zsh{D.3}에서는 열을 선택하거나 원소를 조회했을 때 반환되는 값의 \zsh{dtype}과 결측 표현을 분석했다.

이번에는 이 두 개념을 연결할 것이다.

다음 Series를 생각해 보자.

\begin{dhkcode}
import pandas as pd

s = pd.Series(
    [10, 20, 30],
    index=[2, 0, 1],
)
\end{dhkcode}

여기서 \zsh{s[0]}은 무엇을 반환할까?

\begin{dhksmry}
    - 첫 번째 위치에 저장된 10일까?

    - 레이블 0에 대응하는 20일까?
\end{dhksmry}

정답은 20이다.

\begin{dhkcode}
print(s[0])
print(s.loc[0])
print(s.iloc[0])
\end{dhkcode}

예상 출력:

\begin{dhkzsh}
20
20
10
\end{dhkzsh}

이 차이가 이번 학습의 출발점이다.

Pandas에서 데이터를 선택하는 방법은 단순히 문법의 차이가 아니다.

어떤 기준으로 데이터를 식별하고, 어떤 객체를 반환하며, 그 과정에서 레이블을 어떻게 해석하는가의 차이다.

\subsubsection{이번 학습의 핵심 질문}

이번 \zsh{D.4}에서는 다음 문제를 순서대로 해결할 것이다.

\begin{dhksmry}
1. 레이블(label)과 위치(position)는 정확히 무엇이 다른가?

2. [], .loc, .iloc, .at, .iat은 각각 어떤 객체와 연산을 사용하는가?

3. 하나의 행·열·원소를 선택할 때 반환 타입은 어떻게 결정되는가?

4. 슬라이싱에서 종료 경계가 포함되는 경우와 제외되는 경우는 왜 다른가?

5. 중복 레이블과 불리언 마스크는 선택 결과에 어떤 영향을 주는가?

6. 선택과 할당은 원본 객체의 상태를 어떻게 변화시키는가?
\end{dhksmry}

특히 \zsh{D.4}는 단순한 API 암기가 아니라, Python의 \zsh{__getitem__}, 속성 탐색, 이름 바인딩, 객체 동일성, 슬라이스 객체를 Pandas에 연결하는 학습이다.

\subsubsection{레이블과 위치는 서로 다른 좌표계다}

먼저 앞으로 사용할 공통 에제를 정의한다.

\begin{dhknote}
\begin{dhksmry}[centered]
독립 실행 코드 1 - 두 좌표계의 구성
\end{dhksmry}
\end{dhknote}

\begin{dhkcode}
import pandas as pd

df = pd.DataFrame(
    {
        "text": ["alpha", "beta", "gamma"],
        "label": [1, 0, 1],
        "token_length": [5, 3, 7],
    },
    index=pd.Index(
        [2, 0, 1],
        name="doc_id",
    ),
)

print(df)
\end{dhkcode}

예상 출력:

\begin{dhkzsh}
         text  label  token_length
doc_id
2       alpha      1             5
0        beta      0             3
1       gamma      1             7
\end{dhkzsh}

이 DataFrame은 다음 두 체계로 행을 식별할 수 있다.

\begin{dhkmatrix}[delim=none,style=cards]
title:: 하나의 DataFrame, 서로 다른 행 식별 기준
cols::위치, Index 레이블, text, label
data::[[0,2,alpha,1],[1,0,beta,0],[2,1,gamma,1]]
caption::위치는 현재 배열의 순서이고, Index 레이블은 각 행에 부여된 식별값이다.
\end{dhkmatrix}

첫 번째 행은 위치가 0이지만 레이블은 2다.

두 번째 행은 위치가 1이지만 레이블은 0이다.

따라서 다음 두 표현식은 서로 다른 행을 선택한다.

\begin{dhkcode}
print(df.loc[0, "text"])
print(df.iloc[0, 0])
\end{dhkcode}

예상 출력:

\begin{dhkzsh}
beta
alpha
\end{dhkzsh}

\zsh{df.loc[0, "text"]}는 레이블 0인 행과 레이블 \zsh{"text"}인 열의 교차점을 선택한다.

반면 \zsh{df.iloc[0, 0]}은 첫 번째 행과 첫 번째 열의 교차점을 선택한다.

여기서 중요한 점은 레이블이 정수라는 사실이 그 레이블을 위치로 만들어 주지는 않는다는 것이다.

이 구분은 데이터가 정렬되거나 필터링된 이후 더욱 중요해진다.

행의 위치는 달라질 수 있지만 Index 레이블은 별도의 변경이나 재설정이 없는 한 그대로 유지되기 때문이다.

\subsubsection{\zsh{[]}는 Series와 DataFrame에서 같은 의미일까?}

Python에서 다음 문법은 구독 연산(subscription)이다.

\begin{dhkcode}
obj[key]
\end{dhkcode}

객체 타입에 정의된 \zsh{__getitem__()} 메커니즘을 통해 처리된다.

Pandas의 \zsh{Series}와 \zsh{DataFrame}도 이 문법을 사용하지만, 두 객체가 키를 해석하는 방식은 같지 않다.

\zsh{D.4.3}의 \zsh{df}를 그대로 사용해 보자.

\begin{dhkcode}
s = df["label"]

print(type(s).__name__)
print(s.tolist())

print(type(df[["label"]]).__name__)
print(df[["label"]].shape)

print(df[:2].index.tolist())
\end{dhkcode}

예상 출력:

\begin{dhkzsh}
Series
[1, 0, 1]
DataFrame
(3, 1)
[2, 0]
\end{dhkzsh}

다음 표에서 차이를 확인할 수 있다.

\begin{dhktable}
|표현식|	선택 기준|	반환 타입|
|:---|:---|:---|
|s[0]|	Series의 정수 레이블 0|	스칼라|
|df["label"]|	DataFrame의 열 레이블|	일반적으로 Series|
|df[["label"]]|	열 레이블 목록|	DataFrame|
|df[:2]|	처음 두 행의 위치 범위|	DataFrame|
\end{dhktable}

특히 \zsh{df[:2]}가 행을 선택한다는 점에 주목해야 한다.

\begin{dhkcode}
df["label"]
df[:2]
\end{dhkcode}

두 표현식 모두 \zsh{df[...]} 문법을 사용하지만, 전달된 키의 종류에 따라 선택 대상이 달라지는 것이다.

이것은 \zsh{__getitem__()}이 키의 타입과 형태를 해석해 서로 다른 동작을 수행하기 때문이다.

따라서 DataFrame의 \zsh{[]}를 단순히 열 선택 연산이라고 정의하면 불완전하다.

일반적인 문자열 키는 열 선택에 사용되지만, 슬라이스는 행 선택에 사용될 수 있다.

실전에서는 이 차이로 인한 혼동을 줄이기 위해 행을 선택할 때 \zsh{.loc} 또는 \zsh{.iloc}을 명시하는 편이 좋다.

\subsubsection{\zsh{.loc}과 \zsh{iloc}은 메서드가 아니라 인덱서 객체다}

이번 학습에서 Python 객체 모델과 가장 직접적으로 연결되는 부분이다.

우리는 다음과 같은 문법을 사용한다.

\begin{dhkcode}
df.loc[0, "text"]
df.iloc[0, 0]
\end{dhkcode}

그런데 \zsh{.loc}과 \zsh{.iloc}은 일반적인 메서드가 아니다.

\begin{dhkcode}
print(type(df.loc).__name__)
print(type(df.iloc).__name__)
\end{dhkcode}

예상 출력:

\begin{dhkzsh}
_LocIndexer
_iLocIndexer
\end{dhkzsh}

\zsh{df.loc}은 DataFrame에 연결된 레이블 기반 인덱서(indexer) 객체를 제공하고,

\zsh{df.iloc}은 위치 기반 인덱서 객체를 제공한다.

실행 흐름을 단순화하면 다음과 같다.

\begin{dhkgraph}[uniform]
node n1 (1,1) [span=2] :: DataFrame 객체 df
node n2 (2,1) [accent] :: df.loc//레이블 인덱서
node n3 (3,1) [accent]:: [0, "text"]//"beta"
node n4 (2,2) :: df.iloc//위치 인덱서
node n5 (3,2) :: [0, 0]//"alpha"

split n1 -> n2, n4
edge n2->n3
edge n4->n5
caption::속성 탐색으로 인덱서 객체를 얻고, 그 객체에 구독 연산을 적용하는 흐름이다.
\end{dhkgraph}

이제 다음 표현식을 살펴보자.

\begin{dhkcode}
df.loc[0, "text"]
\end{dhkcode}

Python 문법 수준에서 대괄호 안의 \zsh{0, "text"}는 두 원소를 갖는 튜플로 전달된다.

개념적으로는 다음과 같은 흐름이다.

\begin{dhkcode}
indexer = df.loc
result = indexer.__getitem__((0, "text"))
\end{dhkcode}

이것은 내부 구현을 단순화한 설명이다.

실제 Pandas는 인덱서의 종류와 키 형태에 따라 여러 내부 경로를 사용한다.

반면 다음 표현식은 두 번의 선택 연산을 수행한다.

\begin{dhkcode}
df.loc[0]["text"]
\end{dhkcode}

먼저 \zsh{df.loc[0]}으로 행을 선택한 다음, 반환된 Series에서 \zsh{"text"}를 선택한다.

따라서 두 표현식은 읽기 결과가 같더라도 실행 구조가 다르다.

특히 값을 할당할 때는 이 차이가 중요해진다.

두 단계로 선택한 뒤 값을 변경하는 연쇄 대입(chained assignment)은 Copy-on-Write와 관련된 문제가 있으므로 \zsh{D.7}에서 자세히 다룰 것이다.

\subsubsection{선택 키의 형태에 따라 반환 객체가 달라진다}

\zsh{.loc}과 \zsh{.iloc}은 단순히 데이터를 찾는 기능만 제공하지 않는다.

선택한 축을 유지하는지, 축 하나를 제거하는지에 따라 반환 객체의 차원이 달라진다.

\zsh{D.4.3}의 \zsh{df}를 사용해 보자.

\begin{dhkcode}
a = df.loc[2]
b = df.loc[[2]]
c = df.loc[2, "label"]
d = df.loc[2, ["label"]]
e = df.loc[[2], ["label"]]

print(type(a).__name__, a.shape)
print(type(b).__name__, b.shape)
print(type(c).__name__, c)
print(type(d).__name__, d.shape)
print(type(e).__name__, e.shape)
\end{dhkcode}

예상 출력:

\begin{dhkzsh}
Series (3,)
DataFrame (1, 3)
int64 1
Series (1,)
DataFrame (1, 1)
\end{dhkzsh}

여기서 \zsh{int64}는 \zsh{type(c).__name__}으로 확인한 NumPy 스칼라 타입 이름이다.

\begin{dhktable}
|선택 표현식|	반환 객체|	차원|
|:---|:---|:---|
|df.loc[2]|	Series	|1|
|df.loc[[2]]|	DataFrame	|2|
|df.loc[2, "label"]|	NumPy 정수 스칼라	|0|
|df.loc[2, ["label"]]|	Series	|1|
|df.loc[[2], ["label"]]|	DataFrame	|2|
\end{dhktable}

이 표를 보면 하나의 규칙을 발견할 수 있다.

일반적으로 유일한 레이블을 가진 축에서 스칼라 키 하나로 선택하면 해당 축이 축소되고, 리스트 형태로 선택하면 해당 축이 유지된다.

예를 들어 다음 두 표현식을 비교해 보자.

\begin{dhkcode}
df.loc[2]
df.loc[[2]]
\end{dhkcode}

첫 번째는 행 하나를 Series로 반환한다.

이때 원본 DataFrame의 열 레이블이 반환된 Series의 Index가 된다.

두 번째는 행 하나만 포함하는 DataFrame을 반환한다.

따라서 행 축과 열 축이 모두 유지된다.

다만 이것은 유일한 레이블을 가진 일반적인 경우에 관한 설명이다.

중복 레이블이 있다면 스칼라 레이블을 전달해도 여러 행이나 열이 선택될 수 있다.

\begin{dhknote}
\begin{dhksmry}[centered]
행을 Series로 선택할 때 dtype은 어떻게 될까?
\end{dhksmry}
\end{dhknote}

\zsh{D.4.3}의 DataFrame에는 문자열 열과 정수 열이 함께 존재한다.

\begin{dhkcode}
row = df.loc[2]

print(row.dtype)
\end{dhkcode}

예상 출력:

\begin{dhkzsh}
object
\end{dhkzsh}

하나의 Series는 하나의 dtype을 갖기 때문에, 서로 다른 dtype의 열들을 하나의 행 Series로 구성하면 공통 표현을 위해 dtype이 달라질 수 있다.

이것은 \zsh{D.3}의 자료형 개념과 직접 연결된다.

따라서 특정 열의 dtype을 보존해야 하는 분석에서는 행 전체를 Series로 추출한 결과와 원래 열의 dtype이 반드시 같다고 가정하면 안 된다.

\subsubsection{\zsh{.loc}과 \zsh{.iloc}의 슬라이싱은 왜 종료 경계가 다를까?}

Python의 일반적인 슬라이스는 다음 구조를 사용한다.

\begin{dhkcode}
sequence[start:stop]
\end{dhkcode}

여기서 \zsh{stop}의 위치는 포함되지 않는다.

Pandas의 \zsh{.iloc}도 위치 기반 슬라이싱에서 이 규칙을 따른다.

하지만 \zsh{.loc}은 레이블 기반 슬라이싱에서 일반적으로 종료 레이블을 포함한다.

\begin{dhknote}
\begin{dhksmry}[centered]
독립 실행 코드 2 - 슬라이스 경계 비교
\end{dhksmry}
\end{dhknote}

\begin{dhkcode}
import pandas as pd

df = pd.DataFrame(
    {"score": [90, 80, 70]},
    index=[2, 0, 1],
)

print(df.loc[0:1].index.tolist())
print(df.iloc[0:2].index.tolist())
\end{dhkcode}

예상 출력:

\begin{dhkzsh}
[0, 1]
[2, 0]
\end{dhkzsh}

왜 다를까?

\zsh{df.loc[0:1]}은 레이블 0에서 시작해 레이블 1까지 선택한다.

현재 Index에서 두 레이블이 모두 존재하므로 두 행이 포함된다.

반면 \zsh{df.iloc[0:2]}는 위치 0부터 위치 2 직전까지 선택한다.

따라서 위치0, 1에 해당하는 레이블 2, 0의 행이 반환된다.

이 차이를 다음처럼 정리할 수 있다.

\begin{dhktable}
|구분|	.loc[start:stop]|	.iloc[start:stop]|
|:---|:---|:---|
|기준|	레이블|	정수 위치|
|시작 경계|	포함|	포함|
|종료 경계|	일반적으로 포함|	제외|
|정수 인수의 의미|	정수 레이블|	정수 위치|
|음수 인수|	음수 레이블|	뒤에서부터의 위치|
\end{dhktable}

\begin{dhknote}
\begin{dhksmry}[centered]
두 축을 동시에 슬라이싱하면?
\end{dhksmry}
\end{dhknote}

\begin{dhkcode}
result = df.loc[0:1]
\end{dhkcode}

이 표현식은 행만 선택한다.

반면 DataFrame에서는 행과 열을 동시에 지정할 수 있다.

\begin{dhkcode}
result = df.loc[
    0:1,
    "label":"token_length",
]
\end{dhkcode}

\zsh{D.4.3}의 공통 DataFrame을 사용했다면 결과는 다음과 같다.

\begin{dhkzsh}
        label  token_length
doc_id
0           0             3
1           1             7
\end{dhkzsh}

동일한 영역을 위치 기준으로 선택하면 다음과 같다.

\begin{dhkcode}
result = df.iloc[1:3, 1:3]
\end{dhkcode}

이 경우에는 행 위치 1, 2와 열 위치 1, 2를 선택한다.

여기서 중요한 사실은 \zsh{.loc}의 종료 경계 포함 규칙이 행과 열 모두에 적용된다는 점이다.

또한 \zsh{.loc} 슬라이스는 숫자의 대소관계를 단순 비교하는 연산이 아니다.

Index에 저장된 레이블의 순서와 슬라이싱 가능 조건에 따라 범위를 결정한다.

예를 들어 정렬되지 않은 Index에서 슬라이스 경계 레이블이 존재하지 않으면 \zsh{KeyError}가 발생할 수 있다.

반면 정렬된 Index에서는 존재하지 않는 경계값을 기준으로 범위를 결정할 수 있는 경우도 있다.

따라서 \zsh{.loc} 슬라이싱을 일반적인 숫자 구간 필터링과 동일하게 이해해서는 안 된다.

\subsubsection{\zsh{.at}과 \zsh{.iat}은 무엇이 다른가?}

\zsh{.loc}과 \zsh{.iloc}은 행, 열, 범위, 목록 등을 선택할 수 있는 범용 인덱서다.

반면 \zsh{.at}과 \zsh{.iat}은 하나의 셀에 접근하는 용도로 특화된 인덱서다.

\begin{dhktable}
|인덱서	|기준|	주요 목적|
|:---|:---|:---|
|.loc|	레이블|	행·열·범위 선택|
|.iloc|	위치|	행·열·범위 선택|
|.at|	레이블|	단일 셀 조회·할당|
|.iat|	위치|	단일 셀 조회·할당|
\end{dhktable}

\zsh{D.4.3}의 DataFrame을 사용하면 다음과 같다.

\begin{dhkcode}
print(df.at[2, "label"])
print(df.iat[0, 1])
\end{dhkcode}

예상 출력:

\begin{dhkzsh}
1
1
\end{dhkzsh}

첫 번째 표현식은 행 레이블 2, 열 레이블 \zsh{"label"}을 선택한다.

두 번째 표현식은 첫 번째 행과 두 번째 열을 선택한다.

현재 두 좌표가 동일한 셀을 가리키므로 결과도 같다.

\begin{dhknote}
\begin{dhksmry}[centered]
.at은 언제나 스칼라를 반환할까?
\end{dhksmry}
\end{dhknote}

일반적으로는 그렇지만, 레이블이 중복된 경우에는 예외가 있다.

\begin{dhkcode}
import pandas as pd

df = pd.DataFrame(
    {"score": [10, 20, 30]},
    index=["A", "A", "B"],
)

result = df.at["A", "score"]

print(type(result).__name__)
print(result.tolist())
\end{dhkcode}

Pandas 2.2.3 기준 실행 결과:

\begin{dhkzsh}
Series
[10, 20]
\end{dhkzsh}

왜 스칼라가 아니라 Series가 반환되었을까?

행 레이블 \zsh{"A"}가 두 번 등장하기 때문이다.

\zsh{"A"}라는 레이블만으로는 하나의 행을 유일하게 결정할 수 없다.

이 경우 Pandas의 \zsh{.at}은 단일 셀을 보장할 수 없고, 중복 레이블에 대응하는 여러 값을 반환할 수 있다.

따라서 \zsh{.at}을 사용하는 코드에서 단일 값이 반드시 필요하다면 Index와 열 레이블의 유일성을 먼저 확인하는 것이 좋다.

\begin{dhkcode}
print(df.index.is_unique)
print(df.columns.is_unique)
\end{dhkcode}

예상 출력:

\begin{dhkzsh}
False
True
\end{dhkzsh}

이 문제는 단순한 성능 문제가 아니라 레이블이 데이터를 유일하게 식별하는가에 관한 문제다.

\subsubsection{불리언 마스크를 선택할 때 레이블 정렬이 개입한다}

이 부분은 \zsh{D.5}의 Index alignment와 \zsh{D.6}의 불리언 마스크를 연결하는 핵심 개념이다.

일반적인 NumPy 배열에서 불리언 마스크는 원소의 위치를 기준으로 적용된다.

하지만 Pandas의 \zsh{.loc}에 불리언 Series를 전달하면 Series의 Index 레이블이 선택 과정에 개입할 수 있다.

\begin{dhknote}
\begin{dhksmry}[centered]
독립 실행 코드 3 - 같은 불리언 값, 다른 선택 결과
\end{dhksmry}
\end{dhknote}

\begin{dhkcode}
import pandas as pd

df = pd.DataFrame(
    {"score": [90, 80, 70]},
    index=[2, 0, 1],
)

mask = pd.Series(
    [True, False, True],
    index=[1, 2, 0],
)

by_label = df.loc[mask]
by_position = df.iloc[mask.to_numpy()]

print(by_label.index.tolist())
print(by_position.index.tolist())
\end{dhkcode}

예상 출력:

\begin{dhkzsh}
[0, 1]
[2, 1]
\end{dhkzsh}

두 결과가 왜 다를까?

먼저 \zsh{mask}의 값을 살펴보자.

\begin{dhktable}
|마스크의 레이블	|불리언 값|
|:---|:---|
|1	|True|
|2	|False|
|0	|True|
\end{dhktable}

\zsh{.loc[mask]}는 DataFrame의 Index 레이블에 맞춰 마스크를 해석한다.

따라서 레이블 2는 제외하고, 레이블 0과 1을 선택한다.

반면 다음 표현식은 다르다.

\begin{dhkcode}
mask.to_numpy()
\end{dhkcode}

이 변환을 통해 Index 레이블 정보가 사라지고, 불리언 값의 순서만 남는다.

\begin{dhkzsh}
[True, False, True]
\end{dhkzsh}

이 배열을 \zsh{.iloc}에 전달하면 첫 번째 행과 세 번째 행을 선택한다.

그래서 레이블 2와 1이 반환되는 것이다.

이것은 매우 중요한 차이다.

불리언 Series를 레이블 기준으로 적용하는 것과 불리언 배열을 위치 기준으로 적용하는 것은 서로 다른 연산이다.

특히 다른 DataFrame에서 만든 마스크를 무심코 \zsh{.to_numpy()}로 변환하면 행 순서가 다를 때 잘못된 샘플을 선택할 수 있다.

\begin{dhknote}
\begin{dhksmry}[centered]
레이블이 대응하지 않으면?
\end{dhksmry}
\end{dhknote}

\begin{dhkcode}
bad_mask = pd.Series(
    [True, False],
    index=[2, 0],
)

df.loc[bad_mask]
\end{dhkcode}

현재 DataFrame에는 레이블 1도 존재한다.

그런데 \zsh{bad_mask}에는 레이블 1에 대응하는 값이 없다.

따라서 이 코드는 다음 오류를 발생시킨다.

\begin{dhkzsh}
IndexingError: Unalignable boolean Series provided as indexer ...
\end{dhkzsh}

이 오류는 데이터 길이만의 문제가 아니다.

선택 대상의 Index와 불리언 Series의 Index가 정합하게 대응하지 않는 문제다.

이후 \zsh{D.5}에서는 이러한 정렬이 산술 연산과 할당에서도 어떻게 적용되는지 분석할 것이다.

\subsubsection{선택 실패는 어떤 예외를 발생시키는가?}

레이블 기반 선택과 위치 기반 선택은 실패할 때도 다른 예외를 발생시킬 수 있다.

\zsh{D.4.3}의 DataFrame을 기준으로 다음 코드를 살펴보자.

\begin{dhkcode}
try:
    df.loc[99]
except Exception as e:
    print(type(e).__name__)

try:
    df.iloc[99]
except Exception as e:
    print(type(e).__name__)
\end{dhkcode}

예상 출력:

\begin{dhkzsh}
KeyError
IndexError
\end{dhkzsh}

\zsh{df.loc[99]}는 레이블 99를 찾으려고 한다.

하지만 현재 Index에는 99가 없으므로 \zsh{KeyError}가 발생한다.

반면 \zsh{df.iloc[99]}는 위치 99의 행을 찾으려고 한다.

현재 DataFrame에는 세 개의 행만 있으므로 유효한 위치 범위를 벗어나 \zsh{IndexError}가 발생한다.

\begin{dhknote}
\begin{dhksmry}[centered]
음수 인덱스도 다르게 해석된다
\end{dhksmry}
\end{dhknote}

\begin{dhkcode}
print(df.iloc[-1, 1])
\end{dhkcode}

예상 출력:

\begin{dhkzsh}
1
\end{dhkzsh}

\zsh{.iloc[-1, 1]}은 마지막 행의 두 번째 열을 선택한다.

하지만 다음 표현식은 다르다.

\begin{dhkcode}
df.loc[-1, "label"]
\end{dhkcode}

현재 Index에 -1이라는 레이블이 없으므로 \zsh{KeyError}가 발생한다.

\zsh{.loc}은 음수를 뒤에서부터의 위치로 해석하지 않기 때문이다.

이 차이를 알고 있으면 예외의 종류만 보고도 어떤 선택 기준에서 문제가 발생했는지 추론할 수 있다.

\subsubsection{선택과 할당은 원본 객체를 어떻게 변화시킬까?}

지금까지는 주로 데이터를 읽었다.

하지만 Pandas 인덱서는 값을 변경하는 데에도 사용된다.

다음 예제를 살펴보자.

\begin{dhknote}
\begin{dhksmry}[centered]
독립 실행 코드 4 - DataFrame의 상태 변경
\end{dhksmry}
\end{dhknote}

\begin{dhkcode}
import pandas as pd

original = pd.DataFrame(
    {"score": [90, 80]},
    index=["A", "B"],
)

df = original.copy()
alias = df

df.loc["A", "score"] = 100

print(original.loc["A", "score"])
print(df.loc["A", "score"])
print(alias.loc["A", "score"])

print(df is alias)
\end{dhkcode}

예상 출력:

\begin{dhkzsh}
90
100
100
True
\end{dhkzsh}

이 과정을 Python 객체 모델로 분석해 보자.

첫째, \zsh{original.copy()}가 새로운 DataFrame 객체를 반환한다.

반환 객체가 \zsh{df}라는 이름에 바인딩된다.

둘째, \zsh{alias = df}는 새로운 DataFrame을 생성하지 않는다.

기존 \zsh{df}가 참조하는 객체에 \zsh{alias}라는 이름을 추가로 바인딩한 것이다.

셋째, \zsh{df.loc["A", "score"] = 100}은 기존 DataFrame 객체의 상태를 변경한다.

따라서 같은 객체를 참조하는 \zsh{alias}에서도 변경된 값이 관찰된다.

반면 \zsh{original}은 별도의 DataFrame 객체이므로 원래 값을 유지한다.

여기서 중요한 것은 \zsh{.loc}이 새로운 DataFrame을 반환했느냐가 아니다.

할당문이기 때문에 선택한 위치에 값을 기록하도록 Pandas의 인덱싱 할당 메커니즘이 작동한 것이다.

다만 다음 두 표현식을 동일하게 취급해서는 안 된다.

\begin{dhkcode}
df.loc["A", "score"] = 100
df.loc["A"]["score"] = 100
\end{dhkcode}

첫 번째는 원본 DataFrame에 대한 직접적인 인덱싱 할당이다.

두 번째는 중간에 행 선택 결과를 거치는 연쇄 대입이므로 Copy-on-Write 설정과 Pandas 버전에 따라 원본 변경 여부가 달라질 수 있다.

이 문제는 \zsh{D.7}의 핵심 학습 내용이다.

\subsubsection{ML/Transformer 응용 - 오류 문서 추적과 행 순서 변화}

이번에는 문서 분류 모델의 예측 결과를 분석해 보자.

모델의 예측값과 정답 레이블을 비교해 오류가 발생한 문서를 찾고 모델 신뢰도에 따라 데이터를 정렬한다고 가정한다.

이 과정에서 \zsh{.loc}과 \zsh{.iloc}의 역할이 분명해진다.

\begin{dhknote}
\begin{dhksmry}[centered]
독립 실행 코드 5 - 문서 ID를 기준으로 오류 추적
\end{dhksmry}
\end{dhknote}

\begin{dhkcode}
import pandas as pd

df = pd.DataFrame(
    {
        "text": ["alpha", "beta", "gamma"],
        "label": [1, 0, 1],
        "prediction": [0, 0, 1],
    },
    index=pd.Index(
        ["doc-C", "doc-A", "doc-B"],
        name="doc_id",
    ),
)

scores = pd.Series(
    [0.93, 0.82, 0.64],
    index=["doc-B", "doc-C", "doc-A"],
    name="confidence",
)

df["confidence"] = scores

errors = df.loc[
    df["label"] != df["prediction"],
    ["text", "label", "prediction"],
]

ordered = df.sort_values("confidence")

print(df["confidence"].tolist())
print(errors.index.tolist())
print(errors.loc["doc-C", "text"])

print(ordered.index.tolist())
print(ordered.loc["doc-C", "text"])
print(ordered.iloc[0]["text"])
\end{dhkcode}

예상 출력:

\begin{dhkzsh}
[0.82, 0.64, 0.93]
['doc-C']
alpha
['doc-A', 'doc-C', 'doc-B']
alpha
beta
\end{dhkzsh}

\begin{dhknote}
\begin{dhksmry}[centered]
실행 과정 분석
\end{dhksmry}
\end{dhknote}

먼저 \zsh{scores}의 Index 순서는 다음과 같다.

\begin{dhkzsh}
doc-B, doc-C, doc-A
\end{dhkzsh}

반면 \zsh{df}의 Index 순서는 다음과 같다.

\begin{dhkzsh}
doc-C, doc-A, doc-B
\end{dhkzsh}

그런데 다음 할당을 수행했다.

\begin{dhkcode}
df["confidence"] = scores
\end{dhkcode}

Pandas는 Series를 DataFrame 열에 할당할 때 Index 레이블을 기준으로 값을 대응시킨다.

따라서 \zsh{df["confidence"]}는 현재 DataFrame의 행 순서에 맞춰 다음 값을 갖게 된다.

\begin{dhkzsh}
doc-C → 0.82
doc-A → 0.64
doc-B → 0.93
\end{dhkzsh}

이것은 \zsh{D.5}에서 다룰 Index alignment의 예고편이다.

다음으로 오류 문서를 선택한다.

\begin{dhkcode}
errors = df.loc[
    df["label"] != df["prediction"],
    ["text", "label", "prediction"],
]
\end{dhkcode}

여기서 \zsh{df["label"] != df["prediction"]}은 불리언 Series를 생성한다.

\zsh{.loc}은 그 조건을 만족하는 행과 지정된 열을 선택해 DataFrame을 반환한다.

현재 오류가 발생한 문서는 \zsh{doc-C} 하나다.

마지막으로 신뢰도 기준으로 정렬한다.

\begin{dhkcode}
ordered = df.sort_values("confidence")
\end{dhkcode}

정렬 후 행의 위치는 바뀌었지만 문서 ID는 유지된다.

따라서 다음 표현식은 여전히 \zsh{doc-C}의 텍스트를 반환한다.

\begin{dhkcode}
ordered.loc["doc-C", "text"]
\end{dhkcode}

반면 다음 표현식은 정렬된 DataFrame의 첫 번째 행을 선택한다.

\begin{dhkcode}
ordered.iloc[0]["text"]
\end{dhkcode}

그래서 결과가 \zsh{"beta"}가 되는 것이다.

여기서 얻어야 할 핵심 원칙은 다음과 같다.

샘플의 정체성을 기준으로 데이터를 추적할 때는 안정적인 식별자 레이블을 사용하고, 현재 배열의 순서를 기준으로 처리할 때는 위치 인덱싱을 사용해야 한다.

이 구분을 놓치면 정렬이나 필터링 이후 텍스트와 정답 레이블, 예측 결과가 서로 다른 샘플에 대응하는 문제가 발생할 수 있다.

\subsubsection{핵심 개념 비교 및 판단 기준}

\begin{dhktable}
|인터페이스	|기준|	대표 사용법|	주의할 점|
|:---|:---|:---|:---|
|Series[...]|	키 종류 및 Index에 따른 해석|	s[0]	|정수 레이블과 위치 혼동|
|DataFrame[...]|	주로 열 레이블, 슬라이스 등|	df["label"]	|키 형태에 따라 행·열 선택이 달라짐|
|.loc[...]|	레이블	|df.loc["doc-A"]|	종료 레이블 포함, 중복 레이블|
|.iloc[...]|	정수 위치	|df.iloc[0]	|범위 초과, 슬라이스 종료 위치 제외|
|.at[...]|	레이블	|df.at["doc-A", "label"]|	중복 레이블이면 스칼라 보장 안 됨|
|.iat[...]|	정수 위치	|df.iat[0, 1]|	유효한 단일 위치 필요|
\end{dhktable}

실제 코드를 읽을 때는 다음 판단 순서를 적용하면 좋다.

\begin{dhksmry}
선택 표현식을 분석하는 순서
1. 선택 대상 객체 확인
   - Series인가, DataFrame인가, 인덱서 객체인가?

2. 키의 의미 확인
   - 레이블인가, 정수 위치인가, 슬라이스인가, 목록인가?

3. 축 확인
   - 행을 선택하는가, 열을 선택하는가, 두 축을 동시에 선택하는가?

4. 반환 타입 확인
   - 스칼라인가, Series인가, DataFrame인가?

5. 경계 조건 확인
   - 중복 레이블, 누락된 레이블, 불리언 마스크 정합성은 어떤가?

6. 상태 변화 확인
   - 조회만 하는가, 원본에 할당하는가, 새로운 객체를 이름에 바인딩하는가?
\end{dhksmry}

이 순서를 따르면 대부분의 Pandas 선택 코드를 Python 객체 모델에 근거해 설명할 수 있다.

특히 \zsh{df.loc[...]}과 \zsh{df.iloc[...]}의 차이를 단순히 이름으로 찾기와 번호로 찾기 정도로 암기하지 않는 것이 중요하다.

이 둘은 서로 다른 인덱서 객체이며, 키 해석, 슬라이싱, 불리언 선택, 예외 발생, 데이터 정렬과의 관계까지 달라진다.

\subsubsection{학습 정리와 다음 학습 위치}

이번 \zsh{D.4}에서는 레이블과 위치라는 두 선택 기준을 Python 객체 모델과 연결했다.

핵심은 Pandas의 데이터 선택이 단순한 배열 접근이 아니라, Index 객체와 인덱서 객체가 함께 참여하는 연산이라는 사실이다.

Pandas의 Index는 선택 연산에만 사용되는 것이 아니다.

다음 두 Series를 생각해 보자.

\begin{dhkcode}
import pandas as pd

a = pd.Series(
    [10, 20],
    index=["A", "B"],
)

b = pd.Series(
    [1, 2],
    index=["B", "A"],
)

print(a + b)
\end{dhkcode}

이 연산은 위치 기준으로 \zsh{[11, 22]}를 계산하지 않는다.

레이블을 기준으로 값을 대응시키기 때문에 다음 결과가 발생한다.

\begin{dhkzsh}
A    12
B    21
dtype: int64
\end{dhkzsh}

\zsh{D.5}에서는 이러한 Index alignment가 산술 연산, 비교 연산, DataFrame과 Series의 연산, 결측값 생성, 할당 과정에 어떻게 개입하는지 분석할 것이다.

공식 참고 자료

\begin{dhksmry}
https://pandas.pydata.org/docs/user_guide/indexing.html
https://pandas.pydata.org/docs/user_guide/advanced.html
https://pandas.pydata.org/docs/user_guide/migration.html
\end{dhksmry}
